	Uncertainty poses a major challenge for managers in both the public and private sectors, and addressing it requires the pursuit of information; risk is inversely proportional to one's knowledge of a given subject, and such knowledge is gained by interpreting information obtained through data mining. While traditional forecasting models achieve a certain level of accuracy when dealing with limited datasets, they encounter limitations when higher precision is required or when dealing with complex data—a scenario increasingly common in the era of big data. Accordingly, this study aims to determine how to develop an AI-based predictive model to effectively estimate Brazil's GDP, taking into account the significant differences and similarities among traditional models such as VAR, ARIMA, DFM, and MIDAS. To this end, quantitative methods will be employed to assess the accuracy, computational cost, and interpretability of the models tested, ultimately determining whether AI models can reduce uncertainty.
	
	\keywords{Artificial Intelligence; Brazilian GDP; Quantitative methods; Model testing.}